\documentclass[11pt]{article}
\usepackage[a4paper,margin=25mm]{geometry}
\usepackage{amsmath,amssymb}
\usepackage{hyperref}
\usepackage{url}
\setlength{\emergencystretch}{3em}
\title{Base-four repunits disprove Conjecture 00000000126}
\author{gaochengzhecpu (AI-assisted submission)}
\date{3 October 2026}
\begin{document}
\maketitle

\paragraph{Statement being refuted.}
The conjecture asserts that for every integer $a\ge2$ there are infinitely many primes of the form
$1+a+a^2+\cdots+a^k$.
We show that for $a=4$ the only prime value is $5$.

\paragraph{Definitions and all-length factorization.}
For an integer $a\ge2$ and length $n\ge1$, put
\[
 R_n(a)=1+a+\cdots+a^{n-1}.
\]
Thus $n=k+1$ in the notation of the source. For the empty sum set $R_0(a)=0$.
The finite geometric-sum identity gives
\[
 3R_n(4)+1=4^n=(2^n)^2.
\]
Let $p=2^n$. Since $p$ is not divisible by $3$, write
$p=3q+r$, where $r\in\{1,2\}$. If $n\ge3$, then $p\ge8$.
In the two residue cases we obtain the following explicit integer factorizations:
\[
\begin{array}{c|c|c}
r & R_n(4) & \text{lower bounds on the factors}\\ \hline
1 & q(p+1) & q=(p-1)/3>1,\quad p+1>1\\
2 & (p-1)(q+1) & p-1>1,\quad q+1=(p+1)/3>1 .
\end{array}
\]
Both identities follow from
$3R_n(4)=(p-1)(p+1)$ and cancellation of $3$.
The inequalities use only $p\ge8$. Therefore every $R_n(4)$ with $n\ge3$ is composite.

The remaining lengths give
\[
R_0(4)=0,\qquad R_1(4)=1,\qquad R_2(4)=5.
\]
The first two numbers are not prime and $5$ is prime. Consequently,
\[
 R_n(4)\text{ is prime}\quad\Longleftrightarrow\quad n=2,
\]
and the complete set of prime values in base $4$ is the singleton $\{5\}$.
This contradicts the asserted infinitude for every base.

\newpage
\paragraph{Formal verification and semantic bridge.}
The self-contained Lean 4.19.0 project imports only \texttt{Std}.
The recursive definition \texttt{repunit a n} is the finite sum $R_n(a)$:
the recursion adds $a^n$ on passing from length $n$ to $n+1$.
The definition \texttt{IsPrime p} requires $p\ge2$ and that every factorization
$p=ab$ in natural numbers has $a=1$ or $b=1$; this is the ordinary prime definition.

The theorem \path{geometric_identity} proves the geometric-sum identity by induction for every length.
The theorem \path{pow_two_nonzero_mod_three} proves the two residue cases by induction,
and \path{all_long_repunit_composite} supplies actual natural-number factors
larger than one for every $n\ge3$. These are symbolic proofs over arbitrary natural numbers,
not a bounded search. The two factors are precisely those in the table.

The theorem \path{prime_repunit_exact} proves the equivalence with $n=2$;
\path{every_prime_repunit_is_five} states that every prime value is $5$.
The predicate
\begin{center}\texttt{InfinitelyManyPrimeRepunits}\end{center}
expresses that prime values are unbounded:
for every $B$ there is a positive length giving a prime value larger than $B$.
For a subset of the natural numbers, this is equivalent to infinitude, since a bounded subset is finite.
The theorem \path{base_four_not_infinitely_many} refutes this property at $B=5$,
and \path{conjecture126_counterexample} negates the original universal claim using $a=4$.
In particular, the formal conclusion concerns all possible lengths and distinct prime values.

The project uses no added axioms, proof placeholders, or native decision procedure.
The principal axiom audits contain only Lean's standard \texttt{propext} and \texttt{Quot.sound}.
The Lake configuration treats warnings as errors.
\end{document}
